% Options for packages loaded elsewhere
% Options for packages loaded elsewhere
\PassOptionsToPackage{unicode}{hyperref}
\PassOptionsToPackage{hyphens}{url}
\PassOptionsToPackage{dvipsnames,svgnames,x11names}{xcolor}
%
\documentclass[
  12pt]{article}
\usepackage{xcolor}
\usepackage{amsmath,amssymb}
\setcounter{secnumdepth}{5}
\usepackage{iftex}
\ifPDFTeX
  \usepackage[T1]{fontenc}
  \usepackage[utf8]{inputenc}
  \usepackage{textcomp} % provide euro and other symbols
\else % if luatex or xetex
  \usepackage{unicode-math} % this also loads fontspec
  \defaultfontfeatures{Scale=MatchLowercase}
  \defaultfontfeatures[\rmfamily]{Ligatures=TeX,Scale=1}
\fi
\usepackage{lmodern}
\ifPDFTeX\else
  % xetex/luatex font selection
\fi
% Use upquote if available, for straight quotes in verbatim environments
\IfFileExists{upquote.sty}{\usepackage{upquote}}{}
\IfFileExists{microtype.sty}{% use microtype if available
  \usepackage[]{microtype}
  \UseMicrotypeSet[protrusion]{basicmath} % disable protrusion for tt fonts
}{}
\makeatletter
\@ifundefined{KOMAClassName}{% if non-KOMA class
  \IfFileExists{parskip.sty}{%
    \usepackage{parskip}
  }{% else
    \setlength{\parindent}{0pt}
    \setlength{\parskip}{6pt plus 2pt minus 1pt}}
}{% if KOMA class
  \KOMAoptions{parskip=half}}
\makeatother
% Make \paragraph and \subparagraph free-standing
\makeatletter
\ifx\paragraph\undefined\else
  \let\oldparagraph\paragraph
  \renewcommand{\paragraph}{
    \@ifstar
      \xxxParagraphStar
      \xxxParagraphNoStar
  }
  \newcommand{\xxxParagraphStar}[1]{\oldparagraph*{#1}\mbox{}}
  \newcommand{\xxxParagraphNoStar}[1]{\oldparagraph{#1}\mbox{}}
\fi
\ifx\subparagraph\undefined\else
  \let\oldsubparagraph\subparagraph
  \renewcommand{\subparagraph}{
    \@ifstar
      \xxxSubParagraphStar
      \xxxSubParagraphNoStar
  }
  \newcommand{\xxxSubParagraphStar}[1]{\oldsubparagraph*{#1}\mbox{}}
  \newcommand{\xxxSubParagraphNoStar}[1]{\oldsubparagraph{#1}\mbox{}}
\fi
\makeatother


\usepackage{longtable,booktabs,array}
\newcounter{none} % for unnumbered tables
\usepackage{calc} % for calculating minipage widths
% Correct order of tables after \paragraph or \subparagraph
\usepackage{etoolbox}
\makeatletter
\patchcmd\longtable{\par}{\if@noskipsec\mbox{}\fi\par}{}{}
\makeatother
% Allow footnotes in longtable head/foot
\IfFileExists{footnotehyper.sty}{\usepackage{footnotehyper}}{\usepackage{footnote}}
\makesavenoteenv{longtable}
\usepackage{graphicx}
\makeatletter
\newsavebox\pandoc@box
\newcommand*\pandocbounded[1]{% scales image to fit in text height/width
  \sbox\pandoc@box{#1}%
  \Gscale@div\@tempa{\textheight}{\dimexpr\ht\pandoc@box+\dp\pandoc@box\relax}%
  \Gscale@div\@tempb{\linewidth}{\wd\pandoc@box}%
  \ifdim\@tempb\p@<\@tempa\p@\let\@tempa\@tempb\fi% select the smaller of both
  \ifdim\@tempa\p@<\p@\scalebox{\@tempa}{\usebox\pandoc@box}%
  \else\usebox{\pandoc@box}%
  \fi%
}
% Set default figure placement to htbp
\def\fps@figure{htbp}
\makeatother





\setlength{\emergencystretch}{3em} % prevent overfull lines

\providecommand{\tightlist}{%
  \setlength{\itemsep}{0pt}\setlength{\parskip}{0pt}}



 
\usepackage[]{natbib}
\bibliographystyle{agsm}


\addtolength{\oddsidemargin}{-.5in}%
\addtolength{\evensidemargin}{-1in}%
\addtolength{\textwidth}{1in}%
\addtolength{\textheight}{1.7in}%
\addtolength{\topmargin}{-1in}%
\makeatletter
\@ifpackageloaded{tcolorbox}{}{\usepackage[skins,breakable]{tcolorbox}}
\@ifpackageloaded{fontawesome5}{}{\usepackage{fontawesome5}}
\definecolor{quarto-callout-color}{HTML}{909090}
\definecolor{quarto-callout-note-color}{HTML}{0758E5}
\definecolor{quarto-callout-important-color}{HTML}{CC1914}
\definecolor{quarto-callout-warning-color}{HTML}{EB9113}
\definecolor{quarto-callout-tip-color}{HTML}{00A047}
\definecolor{quarto-callout-caution-color}{HTML}{FC5300}
\definecolor{quarto-callout-color-frame}{HTML}{acacac}
\definecolor{quarto-callout-note-color-frame}{HTML}{4582ec}
\definecolor{quarto-callout-important-color-frame}{HTML}{d9534f}
\definecolor{quarto-callout-warning-color-frame}{HTML}{f0ad4e}
\definecolor{quarto-callout-tip-color-frame}{HTML}{02b875}
\definecolor{quarto-callout-caution-color-frame}{HTML}{fd7e14}
\makeatother
\makeatletter
\@ifpackageloaded{caption}{}{\usepackage{caption}}
\AtBeginDocument{%
\ifdefined\contentsname
  \renewcommand*\contentsname{Table of contents}
\else
  \newcommand\contentsname{Table of contents}
\fi
\ifdefined\listfigurename
  \renewcommand*\listfigurename{List of Figures}
\else
  \newcommand\listfigurename{List of Figures}
\fi
\ifdefined\listtablename
  \renewcommand*\listtablename{List of Tables}
\else
  \newcommand\listtablename{List of Tables}
\fi
\ifdefined\figurename
  \renewcommand*\figurename{Figure}
\else
  \newcommand\figurename{Figure}
\fi
\ifdefined\tablename
  \renewcommand*\tablename{Table}
\else
  \newcommand\tablename{Table}
\fi
}
\@ifpackageloaded{float}{}{\usepackage{float}}
\floatstyle{ruled}
\@ifundefined{c@chapter}{\newfloat{codelisting}{h}{lop}}{\newfloat{codelisting}{h}{lop}[chapter]}
\floatname{codelisting}{Listing}
\newcommand*\listoflistings{\listof{codelisting}{List of Listings}}
\makeatother
\makeatletter
\makeatother
\makeatletter
\@ifpackageloaded{caption}{}{\usepackage{caption}}
\@ifpackageloaded{subcaption}{}{\usepackage{subcaption}}
\makeatother
\usepackage{bookmark}
\IfFileExists{xurl.sty}{\usepackage{xurl}}{} % add URL line breaks if available
\urlstyle{same}
\hypersetup{
  pdftitle={Enabling Experiential Learning: Utah State University's Analytics Solutions Center},
  pdfauthor={Tyler Brough; Polly Conrad; Chris Corcoran; Doug Derrick; Marc Dotson; Kelly Fadel; Chelsea Harding; Sharad Jones; Robert Mills; Reagan Siggard; Brinley Zabriskie},
  pdfkeywords={Experiential Learning, Pedagogy},
  colorlinks=true,
  linkcolor={blue},
  filecolor={Maroon},
  citecolor={Blue},
  urlcolor={Blue},
  pdfcreator={LaTeX via pandoc}}



\begin{document}


\def\spacingset#1{\renewcommand{\baselinestretch}%
{#1}\small\normalsize} \spacingset{1}


%%%%%%%%%%%%%%%%%%%%%%%%%%%%%%%%%%%%%%%%%%%%%%%%%%%%%%%%%%%%%%%%%%%%%%%%%%%%%%

\date{June 15, 2026}
\title{\bf Enabling Experiential Learning: Utah State University's
Analytics Solutions Center}
\author{
Tyler Brough\\
Utah State University\\
and\\Polly Conrad\\
Utah State University\\
and\\Chris Corcoran\\
Utah State University\\
and\\Doug Derrick\\
Utah State University\\
and\\Marc Dotson\\
Utah State University\\
and\\Kelly Fadel\\
Utah State University\\
and\\Chelsea Harding\\
Utah State University\\
and\\Sharad Jones\\
Utah State University\\
and\\Robert Mills\\
Utah State University\\
and\\Reagan Siggard\\
Utah State University\\
and\\Brinley Zabriskie\\
Utah State University\\
}
\maketitle

\bigskip
\bigskip
\begin{abstract}
Higher education is changing. Public trust in the value proposition of
universities and colleges has eroded in recent decades while their
traditional institutional function of storing, creating, and
disseminating knowledge has been challenged by content delivery on the
internet and private industry investment in research and development. In
response, there has been a growing emphasis within higher education on
what university and colleges can uniquely provide students, including
experiential learning. However, there is a dearth of research on how to
enable experiential learning effectively. We detail how the Analytics
Solutions Center at Utah State University has incentivized alignment of
its faculty and industry partners to enable transformative experiential
learning experiences for students.
\end{abstract}

\noindent%
{\it Keywords:} Experiential Learning, Pedagogy
\vfill

\newpage
\spacingset{1.9} % DON'T change the spacing!


\section{Introduction}\label{sec-intro}

Higher education is changing. Public trust in the value proposition of
universities and colleges has eroded in recent decades while their
traditional institutional function of storing, creating, and
disseminating knowledge has been challenged by content delivery on the
internet and private industry investment in research and development. In
response, there has been a growing emphasis within higher education on
what university and colleges can uniquely provide students, including
experiential learning.

Experiential learning is quite simply learning by doing. It includes
real-world experience, internships, and project-based learning. While
there is a growing body of research on its benefits, there is a dearth
of research on how to enable experiential learning effectively. In this
paper, we detail how the Analytics Solutions Center in the Jon M.
Hunstman School of Business at Utah State University has incentivized
alignment of its faculty and industry partners to enable transformative
experiential learning experiences for students.

\section{Experiential Learning}\label{sec-literature}

\begin{tcolorbox}[enhanced jigsaw, arc=.35mm, bottomrule=.15mm, bottomtitle=1mm, breakable, colback=white, colbacktitle=quarto-callout-note-color!10!white, colframe=quarto-callout-note-color-frame, coltitle=black, left=2mm, leftrule=.75mm, opacityback=0, opacitybacktitle=0.6, rightrule=.15mm, title=\textcolor{quarto-callout-note-color}{\faInfo}\hspace{0.5em}{Note}, titlerule=0mm, toprule=.15mm, toptitle=1mm]

\textbf{Reagan} with Polly, Chelsea, Marc, Sharad, Bob, Brinley?

\begin{itemize}
\tightlist
\item
  Literature review on the benefits of experiential learning
\item
  Reference to relevant pedagogical frameworks
\item
  Literature (if any?) on how to enable experiential learning
  effectively
\item
  Discussion of the barriers in higher education to aligning incentives
\end{itemize}

\end{tcolorbox}

\section{The Analytics Solutions Center}\label{sec-asc}

The Analytics Solutions Center (ASC) in the Jon M. Hunstman School of
Business at Utah State University, Utah's land-grant institution, is a
faculty-supervised, industry-partnered applied learning environment.
Undergraduate and graduate students in the ASC complete data work on
authentic client-facing projects in data analytics, information systems,
data engineering, artificial intelligence, and related STEM and
technology domains. Student teams collaborate with industry, nonprofit,
public-sector, and university partners on interdisciplinary projects
spanning analytics, computing, public-sector, and community-focused STEM
applications. The ASC operates as a layered mentoring ecosystem: faculty
mentors provide disciplinary and methodological guidance, advanced
graduate and undergraduate students serve as near-peer mentors within
project teams, and industry clients engage directly with student teams
throughout the project lifecycle. Sponsoring organizations (including
regional companies, national corporations, nonprofits, and public-sector
entities) develop project scopes, provide data and domain access, and in
many cases provide financial sponsorship for student work.

Unlike traditional course-based project experiences, the ASC functions
as a continuous, institutionally embedded learning environment in which
students may engage across multiple semesters while participating in
increasingly advanced mentoring and technical roles. Students
participate in recurring team meetings, technical mentorship structures,
client presentations, collaborative code and data workflows, and peer
mentoring roles that evolve as students gain experience within the
center. Additionally, the ASC maintains relationships with an ecosystem
of student clubs and professional organizations across campus, including
chapters of the Association for Information Systems and PyData, to
provide students with additional opportunities for technical and
professional development. The ASC's focus on applied learning via
mentored projects while supporting and being supported by special
interest and professional activities creates a community of practice.

Since its founding in 2022, the ASC has engaged over 180 students across
more than 100 projects, with participants representing 34 academic
disciplines across STEM, computing, business, and social science fields.
This wide-ranged engagement reflects Utah State's land-grant mission to
serve a broad, interdisciplinary student population. The center has
developed institutional partnerships at local, state, and national
levels, including collaborations with Utah State Extension and statewide
campuses, Utah State's Office of Data Analytics, and public-sector
organizations across Utah. Notable aspects of the ASC include:

\begin{itemize}
\tightlist
\item
  A project partnership with international Meta Platforms, Inc
\item
  Ongoing data engineering and causal modeling work with the local Cache
  County School District
\item
  A Utah Legislature-funded data engineering and analytics project
  (HB0249, \$300K) focused on state-level fiscal risk modeling and
  decision support
\item
  Recognition in the official state legislature audit for the ASC's
  value-generating research, sustained external demand and institutional
  support through continued partnerships and sponsored project funding,
  and 98\% career/graduate program placement rate
\end{itemize}

The ASC model has already demonstrated replicability. The University of
Rhode Island is currently working with ASC members to design a
comparable center and has committed \$500,000 toward making it
operational. This unsolicited replication effort is evidence that the
model addresses a recognized gap at peer institutions. Land-grant
universities across the country increasingly seek to integrate applied
analytics into their curricula but lack a studied, documented framework
for doing so. The ASC's success and the possibility to replicate it at
other institutions requires an alignment of incentives between faculty,
industry partners, and university leadership. These three areas of
alignment are consistent with the ASC's three-fold mission: empowering
students, providing value to external partners, and developing
cross-disciplinary teams to address complex problems.

First, to empower students through immersive, hands-on, real-world
challenges and experiences requires faculty to be incentivized to engage
as mentors. ASC faculty have their mentorship role included as part of
their role statement. These mentors are both professional practice
faculty and tenure-track faculty where project mentorship serves as a
course release and may dovetail into academic research. The formal
inclusion of mentorship responsibilities in faculty role statements is
crucial for the ASC to operate. Faculty participation isn't a service,
it's a core responsibility, and the success of ASC students is a key
component of ASC faculty promotion and retention.

Second, to provide value to external partners and serve as an economic
driver for the region requires industry partners to be incentivized to
engage as both clients and sponsors. With every ASC project, industry
partners clearly understand that they are primarily providing students
with a real-world learning experience and secondarily receiving value
from the work the students accomplish. A key part of providing that
real-world learning experience while ensuring that the students provide
work that is of value to the partner is a regular meeting with the
students and the mentor to ensure the project remains on track and
within scope. The ASC thus operates on donation-based funding rather
than serving a student consultancy. This model is a crucial part of how
the ASC operates and keep the focus squarely on student learning.

Third, to develop cross-disciplinary teams that address complex
analytical problems and create technological innovations requires
university leadership to be incentivized to help fund the center and
support its operations. The ASC has been supported by the Jon M.
Hunstman School of Business, specifically by the Department of Data
Analytics and Information Systems, which has provided support for the
center and serves as the academic home for the cross-disciplinary
assortment of professional practice and tenure-track ASC faculty
mentors. This support is crucial for the ASC to operate and have the
necessary resources for students and faculty mentors to succeed.

\begin{tcolorbox}[enhanced jigsaw, arc=.35mm, bottomrule=.15mm, bottomtitle=1mm, breakable, colback=white, colbacktitle=quarto-callout-note-color!10!white, colframe=quarto-callout-note-color-frame, coltitle=black, left=2mm, leftrule=.75mm, opacityback=0, opacitybacktitle=0.6, rightrule=.15mm, title=\textcolor{quarto-callout-note-color}{\faInfo}\hspace{0.5em}{Note}, titlerule=0mm, toprule=.15mm, toptitle=1mm]

\textbf{Marc} with Tyler, Polly, Chris, Doug, Kelly, Sharad?

\begin{itemize}
\tightlist
\item
  Reference ASC recognition and success (number of sponsors, money
  raised, placement numbers, etc.)
\item
  Include additional anecdotal evidence of sucess (e.g., student and
  partner testimonials)
\end{itemize}

\end{tcolorbox}

\section{Conclusion}\label{sec-conclusion}

\begin{tcolorbox}[enhanced jigsaw, arc=.35mm, bottomrule=.15mm, bottomtitle=1mm, breakable, colback=white, colbacktitle=quarto-callout-note-color!10!white, colframe=quarto-callout-note-color-frame, coltitle=black, left=2mm, leftrule=.75mm, opacityback=0, opacitybacktitle=0.6, rightrule=.15mm, title=\textcolor{quarto-callout-note-color}{\faInfo}\hspace{0.5em}{Note}, titlerule=0mm, toprule=.15mm, toptitle=1mm]

\begin{itemize}
\tightlist
\item
  Summarize the benefits of and challenges for experiential learning in
  higher education
\item
  Summarize what the ASC has done to address these challenges and enable
  experiential learning
\item
  Plans for future research about the ASC, including tracking any
  additional student outcomes as part of ESL/NSF grant?
\item
  Reference the vision to be ``the premier center in the Intermountain
  West for applied academics at the forefront of data analytics and
  technological innovation.''
\item
  Discuss implications for Utah State University, including Sage
  Academy, specifically and higher education more broadly
\end{itemize}

\end{tcolorbox}


\bibliography{references.bib}



\end{document}
